% (User input window)
% (Source tag) /killo-golden
% Timestamp: 2026-02-18 15:11:15 (Asia/Singapore)
%
% Rewritten using the existing notation and auditable interfaces of this manuscript, without introducing manual numbering.

\begin{conclusion}[Planar-cubic structure of the cubic resolvent curve and ramification type over the fiber at infinity \texorpdfstring{$(2,1)$}{(2,1)}]\label{con:xi-terminal-zm-resolvent-cubic-plane-cubic-geometry}
Adopt the cubic resolvent from Conclusion \ref{con:xi-terminal-zm-generic-galois-s4}:
\[
\mathscr R(z,y)
:=z^3+(1+2y)z^2-(1+4y+4y^2)z-(1+5y+13y^2+8y^3)\in\ZZ[z,y],
\]
and let \(\mathcal R\subset \mathbb A^2_{z,y}\) be the affine curve \(\mathscr R(z,y)=0\), with homogeneous closure \(\overline{\mathcal R}\subset \mathbb P^2_{[Z:Y:W]}\) (\(z=Z/W,\ y=Y/W\)).
Then:
\begin{enumerate}
  \item \(\overline{\mathcal R}\) is a nonsingular plane cubic. Hence \(\overline{\mathcal R}\) itself is an elliptic curve over \(\QQ\) (no further normalization is required).
  \item The projection
  \[
  \pi^{(3)}_y:\ \overline{\mathcal R}\longrightarrow \mathbb P^1_y,\qquad [Z:Y:W]\longmapsto [Y:W]
  \]
  is a degree-\(3\) covering. Its branch image is exactly characterized by the discriminant: viewing \(\mathscr R\) as a cubic polynomial in \(z\),
  \[
  \mathrm{Disc}_{z}\bigl(\mathscr R(\cdot,y)\bigr)=-y(y-1)P_{\mathrm{LY}}(y),
  \]
  so simple ramification (index \(2\)) occurs at \(y=0\), \(y=1\), and the three points defined by \(P_{\mathrm{LY}}(y)=0\).
  \item Over the fiber at infinity \(y=\infty\), the fiber of \(\pi^{(3)}_y\) consists of two points at infinity, with pole-order distribution \(2+1\) (ramification type \((2,1)\)).
  More precisely, the intersection of the line \(W=0\) with \(\overline{\mathcal R}\) is
  \[
  [2:1:0]\ \text{(intersection multiplicity \(1\))},\qquad [-2:1:0]\ \text{(intersection multiplicity \(2\))},
  \]
  hence \([-2:1:0]\) is the unique ramification point above \(y=\infty\) (ramification index \(2\)), while \([2:1:0]\) is unramified.
\end{enumerate}
\end{conclusion}
\begin{proof}[Proof sketch]
Take the homogenization
\[
\mathscr R_h(Z,Y,W)
=Z^3+(W+2Y)Z^2-(W^2+4YW+4Y^2)Z-(W^3+5YW^2+13Y^2W+8Y^3).
\]
On the line at infinity \(W=0\), one has
\[
\mathscr R_h(Z,Y,0)=Z^3+2YZ^2-4Y^2Z-8Y^3=(Z-2Y)(Z+2Y)^2,
\]
yielding the two points \([2:1:0]\), \([-2:1:0]\) with multiplicities \(1,2\).
Since \(y=Y/W\), these multiplicities are exactly the pole orders of \(y\) at the corresponding points.
Therefore the fiber type over \(y=\infty\) is \(2+1\), and only \([-2:1:0]\) has ramification index \(2\).
\par
As a cubic in \(z\), its discriminant identity has already been established in Conclusion \ref{con:xi-terminal-zm-generic-galois-s4}; thus the finite branch-value set and ramification indices follow immediately.
Finally, smoothness of \(\overline{\mathcal R}\) can be verified on the affine chart \(W=1\) and the chart at infinity \(Y=1\) by checking that the gradient never vanishes simultaneously.
Equivalently, for the projective curve \(\mathscr R_h=0\), the system \(\mathscr R_h=\partial_Z\mathscr R_h=\partial_Y\mathscr R_h=\partial_W\mathscr R_h=0\) has no solution over an algebraic closure.
This completes the proof.
\end{proof}

